The software ships with a CMake environment, which will configure and build the
programmes. It is recommended to configure and build the executables in a
separate build directory. This also allows to have several builds with different
options from the same source code directory.

\subsection{Prerequisites}

In order to compile the programmes the {\ttfamily LAPACK}~\cite{lapack:web}
library (fortran version) needs to be installed. CMake will search for the
library in all default directories. Also the latest version (tested is version
1.2.3) of {\ttfamily C-LIME}~\cite{lime:web} must be available, which is used as
a packaging scheme to read and write gauge configurations and propagators to
files.

\subsection{Configuring the hmc package}
\label{sec:config}

The build system uses CMake to configure and build the hmc package. The
following list gives all options (OFF by default unless specified):
\begin{itemize}
\item {\ttfamily CMAKE\_POSITION\_INDEPENDENT\_CODE}: Build a position independent
  code. ON by default.
\item {\ttfamily BUILD\_SHARED\_LIBS}: Build the shared version of the hmc library.
\item {\ttfamily TM\_USE\_FFTW}: Enable fftw support. 
\item {\ttfamily TM\_USE\_CUDA}: Enable CUDA support.
\item {\ttfamily TM\_USE\_HIP}: Enable HIP support (AMD or NVidia GPUs)
\item {\ttfamily TM\_USE\_DDalphaAMG}: Enable DDalphaAMG support.
\item {\ttfamily TM\_USE\_LEMON}: Use the lemon io library.
\item {\ttfamily TM\_USE\_OMP}: Enable OpenMP ({\bf ON} by default)
\item {\ttfamily TM\_FIXEDVOLUME}: Fix volume at compile time.
\item {\ttfamily TM\_ENABLE\_ALIGNMENT}: Automatically or expliclty align arrays to
  byte number. auto, none, 16, 32, 64.
\item {\ttfamily TM\_USE\_GAUGE\_COPY}: Enable use of a copy of the gauge field (ON
  by default). See section \ref{sec:dirac} for details on this option. It will
  increase the memory requirement of the code.
\item {\ttfamily TM\_USE\_HALFSPINOR}: Use a Dirac Op. with halfspinor exchange (ON
  by default). See sub-section \ref{sec:dirac} for details. 
\item {\ttfamily TM\_USE\_QUDA}: Enable QUDA support.
\item {\ttfamily TM\_USE\_SHMEM}: Use shmem API.
\item {\ttfamily TM\_ENABLE\_WARNINGS}: Enable all warnings (ON by default).
\item {\ttfamily TM\_ENABLE\_TESTS}: Enable tests.
\item {\ttfamily TM\_USE\_QPHIX}: Enable QPhiX.
  \begin{itemize}
  \item {\ttfamily TM\_QPHIX\_SOALEN}: QPhiX specific parameter (default is 4)
  \item \textcolor{red}{{\ttfamily QPHIX\_DIR}}: Directory where QPhiX is installed.
    The QPhiX current CMake build system does not export all information (
    include and lib directories) that are needed to compile hmc.
  \item \textcolor{red}{\ttfamily QMP\_DIR}: Directory where QMP is installed (
    QPhiX dependency).
    The QPhiX current CMake build system does not export all information about the
    include and lib directories nor its dependencies (QMP in that case).
  \end{itemize}
\item {\ttfamily TM\_USE\_MPI}: Enable MPI support.
  \begin{itemize}
  \item {\ttfamily TM\_PERSISTENT\_MPI}: Use persistent MPI calls for halfspinor.
  \item {\ttfamily TM\_NONBLOCKING\_MPI}: Use non-blocking MPI calls for spinor and
    gauge.
  \item {\ttfamily TM\_MPI\_DIMENSION}: Use $n$ dimensional parallelisation ($XYZT$)
    [default=4]. The number of parallel directions can be specified. $1, 2, 3$ and $4$
    dimensional parallelisation is supported.
  \item {\ttfamily TM\_USE\_LEMON} Use the lemon io library
  \end{itemize}
\end{itemize}

The following minimal list of commands will configure and build the hmc package with
minimal dependencies
\begin{verbatim}
mkdir build
cd build
cmake -DCMAKE_INSTALL_PREFIX=/my_path -DCMAKE_PREFIX_PATH=/my_c_line_path ..
make -j
make install
\end{verbatim}

These instructions assume that the {\ttfamily c-lime} package is installed in {\ttfamily
  /my\_c\_line\_path}. By default {\ttfamily CMAKE\_PREFIX\_PATH} variable is a list
of paths separated by a semi-colunm containing the path of all installed to
dependencies.

Adding {\ttfamily -DTM\_USE\_MPI=ON} will enable MPI support with parallelization
over spatial and temporal dimensions. The command line is then
\begin{verbatim}
cmake -DCMAKE_INSTALL_PREFIX=/my_path -DCMAKE_PREFIX_PATH=/my_c_line_path -DTM_USE_MPI=ON ..
\end{verbatim}
We can combine it with the lemon-io library (isntalled in {\ttfamily /my\_lemon\_path})
\begin{verbatim}
cmake -DCMAKE_INSTALL_PREFIX=/my_path \
      -DCMAKE_PREFIX_PATH="/my_c_line_path;/my_lemon_path" \
      -DTM_USE_MPI=ON \
      -DTM_USE_LEMON=ON ..
\end{verbatim}

{\ttfamily QUDA} support (installed in {\ttfamily my\_quda\_path}) can be added with
\begin{verbatim}
cmake -DCMAKE_INSTALL_PREFIX=/my_path \
      -DCMAKE_PREFIX_PATH="/my_c_line_path;/my_lemon_path;\my_quda_path" \
      -DTM_USE_MPI=ON \
      -DTM_USE_LEMON=ON \
      -DTM_USE_QUDA \
      -DTM_USE_CUDA=ON \
      -DCMAKE_CUDA_ARCHITECTURES=90 ..
\end{verbatim}
Note that the command assumes that QUDA is compiled with CUDA support. AMD GPU
are also supported after replacing {\ttfamily -DTM\_USE\_CUDA=ON} with
{\ttfamily -DTM\_USE\_HIP=ON} and compiling {\ttfamily QUDA} with {\ttfamily
  HIP} support. The {\ttfamily ROCM} architecture is defined by the variable
{\ttfamily CMAKE\_HIP\_ARCHITECTURES=gfxxxx}.

{\ttfamily QPhiX} and/or {\ttfamily DDalphaAMG} support can be added with
\begin{verbatim}
cmake -DCMAKE_INSTALL_PREFIX=/my_path \
      -DCMAKE_PREFIX_PATH="/my_c_line_path;/my_lemon_path;/my_quda_path;/my_path_ddalphaamg" \
      -DTM_USE_MPI=ON \
      -DTM_USE_LEMON=ON \
      -DTM_USE_QUDA=ON \
      -DTM_USE_CUDA=ON \
      -DCMAKE_CUDA_ARCHITECTURES=90 \
      -DTM_USE_QPHIX=ON \
      -DQPHIX_DIR=/my_qphix_dir \
      -DTM_USE_DDalphaAMG=ON \
      -DQMP_DIR=/my_qmp_dir \
      -DTM_USE_OMP=ON ..
\end{verbatim}
{\ttfamily QPhiX} cmake config support is incomplete and requires both the {\ttfamily QPhiX}
and {\ttfamily QMP} installation directories to work properly.

CMake has several relevant specific options that control the build. Compiler
options are defined by the variable {\ttfamily CMAKE\_C\_FLAGS} and {\ttfamily
  CMAKE\_CXX\_FLAGS}. CUDA and HIP compilations options are controlled by their
equivalent {\ttfamily CMAKE\_\{CUDA/HIP\}\_FLAGS}. 

Adding for instance {\ttfamily -GNinja} to the {\ttfamily CMake} command line will use
{\ttfamily ninja} instead of {\ttfamily make}.

%%% Local Variables: 
%%% mode: latex
%%% TeX-master: "main"
%%% End: 
